\section{From a nonzero algebraic norm to a complete rational intersection list}
\label{sec:norm-enumeration}

The earlier norm packets give sharp bounds on the number of distinct
points in a supplied list \cite{Demand,Coherent}. For an \emph{explicitly supplied section},
their finite data yield more: a complete rational-point list, with a
primitive recursive coordinate bound.  The additional ingredient is
the elementary rational-root argument.  It is useful to state this
separately, because a list of actual sections and the assertion that
every tuple admits some sections have different computational content.

\subsection{A total bounded rational-root procedure}

Encode a rational as a signed numerator and a positive denominator in
lowest terms, and a rational polynomial as a finite coefficient list.
The zero polynomial is detected by a bounded coefficient scan.
For a nonzero polynomial \(H\in\mathbb Q[T]\), multiply by the positive
product of all coefficient denominators.  This gives a nonzero integer
polynomial \(h\), with precisely the same rational zeros.  Scan for its
first and last nonzero coefficients and write
\[
 h(T)=T^e(c_0+c_1T+\cdots+c_dT^d),
 \qquad c_0c_d\ne0.
\]
Here \(d\) may be zero.  Put
\[
 B(H)=\max\{1,|c_0|,|c_d|\}.
 \tag{RR1}\label{eq:norm-root-bound}
\]
Define \(\operatorname{Roots}_{\mathbb Q}(H)\) by scanning all pairs
\[
 -B(H)\le p\le B(H),\qquad 1\le q\le B(H),\qquad
 \gcd(|p|,q)=1,
\]
and retaining \(p/q\) exactly when \(H(p/q)=0\).  Each retained
rational has a unique representation in this scan.  The totalized
procedure returns a failure tag on the zero polynomial; it does not
return an allegedly complete finite list in that case.

\begin{lemma}[Bounded rational roots]\label{lem:norm-rational-roots}
For every nonzero rational polynomial \(H\), the displayed procedure
returns exactly its rational roots, without repetition.  Every root
has reduced numerator and denominator bounded in absolute value by
\(B(H)\).  The procedure and its correctness proof use only primitive
recursion and quantifier-free induction on finite arithmetic data.
\end{lemma}

\begin{proof}
Soundness follows from the final exact substitution test.  Let
\(p/q\), with \(q>0\) and \(\gcd(|p|,q)=1\), be a root.
If \(p=0\), reducedness gives \(q=1\), so the scan includes it.
Suppose \(p\ne0\).  Cancelling the nonzero factor \((p/q)^e\)
and multiplying by \(q^d\) gives
\[
 c_0q^d+c_1pq^{d-1}+\cdots+c_dp^d=0.
 \tag{RR2}\label{eq:norm-cleared-root}
\]
For \(d=0\), this contradicts \(c_0\ne0\), so there is no such
root.  For \(d>0\), reduction modulo \(p\) shows
\(p\mid c_0q^d\), and reduction modulo \(q\) shows
\(q\mid c_dp^d\).  Euclid's lemma, applied to the coprime pairs
\((p,q^d)\) and \((q,p^d)\), gives
\(p\mid c_0\) and \(q\mid c_d\).  The coefficients are nonzero,
so \(|p|\le|c_0|\) and \(q\le|c_d|\).  Thus every root occurs in
the scan.

For an elementary arithmetic implementation of the cancellation step,
Bezout coefficients for \(p,q\) are computed by a Euclidean algorithm
with fuel bounded by \(|p|+q+1\).  Multiplying the Bezout relation
finitely many times gives the required coprimality of powers.
All integer divisions, signs and zero tests are decidable.  There are
at most \((2B(H)+1)B(H)\) candidate pairs.  Degree, the coefficient
identities, and the scan prefix are finite arrays; correctness is
therefore proved by Booleanized bounded invariants.  No assertion
about an infinite set of roots is used as an induction formula.
\end{proof}

The bound is deliberately coarse.  Divisor enumeration, content
removal and synthetic division can substantially reduce the computation,
but are unnecessary for the primitive recursive theorem.  A nonzero
constant polynomial gives the empty list.  A polynomial \(cT^e\),
with \(e>0\), gives exactly \(0\).  Thus neither zero constant terms
nor constant quotients require an exception left outside the algorithm.

\subsection{The norm and the intersection algorithm}

Let
\[
 F(X,Y)=Y^m+f_{m-1}(X)Y^{m-1}+\cdots+f_0(X),
 \qquad m\ge1,
\]
with \(f_j\in\mathbb Q[X]\), and let \(G\in\mathbb Q[X,Y]\)
be arbitrary.  Monic division in \(Y\) identifies the quotient by
\(F\) with a free \(\mathbb Q[X]\)-module generated by
\(1,Y,\ldots,Y^{m-1}\).  Reduce each \(GY^j\) modulo \(F\),
and let \(M_G(X)\) be the matrix whose \(j\)-th column is that
coefficient vector.  Define the polynomial norm
\[
 N_G(X)=\det M_G(X)\in\mathbb Q[X].
\]
Reduction and the determinant can be computed by finite coefficient
operations.  In particular, scanning the \(m^m\) length-\(m\)
index tuples and retaining the permutations is an explicit primitive
recursive determinant implementation.  No general algorithm for ideal
membership, algebraic closure or factorization is needed.

\begin{lemma}[Common zeros give norm zeros]
\label{lem:norm-intersection-vanishing}
If \((a,b)\in\mathbb Q^2\) satisfies \(F(a,b)=G(a,b)=0\), then
\(N_G(a)=0\).
\end{lemma}

\begin{proof}
Put \(v=(1,b,\ldots,b^{m-1})\).  Evaluating each monic reduction
identity at \((a,b)\) gives
\[
 vM_G(a)=G(a,b)v=0.
\]
The adjugate identity, proved by finite determinant expansion, yields
\(vM_G(a)\operatorname{adj}(M_G(a))=N_G(a)v\).  Its first
coordinate is \(N_G(a)\), because \(v_0=1\).  It therefore vanishes.
This proof allows repeated roots, singular fibers and \(m=1\).
\end{proof}

Assume now that the finite coefficient test verifies \(N_G\ne0\).
First compute the finite list
\[
 A=\operatorname{Roots}_{\mathbb Q}(N_G).
\]
For each \(a\in A\), form \(F_a(Y)=F(a,Y)\), which is monic
of degree \(m\), hence nonzero.  Compute
\(B_a=\operatorname{Roots}_{\mathbb Q}(F_a)\).  The output list is
\[
 \mathcal I(F,G)
 =\bigl[(a,b):a\in A,\ b\in B_a,\ G(a,b)=0\bigr].
 \tag{NI}\label{eq:norm-intersection-list}
\]
All tests are exact rational tests.  A norm root can be extraneous;
the last test is essential.  The empty list is a permitted result.

\begin{theorem}[Complete rational intersection enumeration]
\label{thm:norm-complete-intersection}
For every monic \(F\) as above and every \(G\) with \(N_G\ne0\),
\(\mathcal I(F,G)\) is a duplicate-free complete list of the
common rational zeros of \(F,G\).  It is produced by one primitive
recursive algorithm, uniformly in the finite coefficient arrays and
the numerical degree \(m\).

An explicit primitive recursive height bound is
\[
 H(F,G)=\max\Bigl(\{1,B(N_G)\}
       \cup\{B(F_a):a\in A\}\Bigr).
 \tag{NH}\label{eq:norm-intersection-height}
\]
Here the height of a rational pair is the maximum of the absolute
values of its reduced numerators and its positive denominators.
Every common rational zero has height at most \(H(F,G)\).
\end{theorem}

\begin{proof}
If a pair is output, its first coordinate is in \(A\), its second
coordinate is an exact root of \(F_a\), and the final test gives
\(G(a,b)=0\).  Thus it is a common zero.  Conversely, any common
zero \((a,b)\) satisfies \(N_G(a)=0\) by
Lemma~\ref{lem:norm-intersection-vanishing}; completeness of the first
root scan puts \(a\) in \(A\).  Since \(F_a(b)=0\), completeness
of the second scan puts \(b\) in \(B_a\).  The final test passes.
Each root scan uses canonical fractions without repetitions, and
the first coordinates distinguish the blocks, so the pair list has
no repetitions.  The height assertion follows directly from
Lemma~\ref{lem:norm-rational-roots} in the two scans.

The number of first-coordinate candidates is bounded by
\((2B(N_G)+1)B(N_G)\).  After this finite list is computed, the
maximum of its specialized bounds is another primitive recursive
fold.  These numerical bounds govern the second scans.  Computing a
bound from a previously computed finite list is ordinary composition
and primitive recursion; it is not unbounded minimization.  Polynomial
evaluation and rational normalization are fixed primitive recursive
programs on coefficient and point codes.  Soundness and completeness
are obtained from bounded scan invariants and the displayed finite
identities.  Thus the written argument supplies an iPRA proof strategy
without induction over an unbounded root predicate.
\end{proof}

The requirement \(N_G\ne0\) is a sufficient certificate, not a
necessary characterization of finite rational intersections.
For example, a zero norm can arise from a common algebraic component
with no rational points.  The algorithm reports an unsupported input
when the norm is identically zero; it draws no conclusion that the
rational intersection is infinite.  In particular \(G=0\) is rejected.
No irreducibility or nonsingularity hypothesis is needed on supported
inputs.  The integer points are obtained simply by retaining output
coordinates whose denominators are one.

\subsection{The norm also constructs a finite multiplication certificate}
\label{sec:norm-to-module}

The norm interface can be connected directly to the finite-module
enumeration theorem.  It supplies an explicit certificate constructor
for the entire supported monic-plane class, not merely an independent
root-finding algorithm.

\begin{proposition}[A norm-to-module certificate constructor]
\label{prop:norm-to-module}
Given \(F,G\) as above with nonzero norm of degree \(d>0\), a
primitive recursive procedure constructs a unit-and-action certificate
for the ideal \((F,G)\), using the \(md\) monomials
\[
             X^iY^j\qquad(0\le i<d,\ 0\le j<m).
\]
If \(N_G\) is a nonzero constant, it instead constructs an ideal
identity for \(1\), certifying that there are no common zeros.
\end{proposition}

\begin{proof}
Retain the quotients from the monic reductions used to construct
\(M_G\): for \(0\le j<m\),
\[
 GY^j=\sum_{i<m}(M_G)_{ij}Y^i+Fh_j.
\]
For the rows \(b=(1,Y,\ldots,Y^{m-1})\) and
\(h=(h_0,\ldots,h_{m-1})\), this says
\(bM_G=Gb-Fh\).  Right multiplication by the adjugate gives
\[
 N_G b=G b\operatorname{adj}(M_G)
           -Fh\operatorname{adj}(M_G).
\]
Its first coordinate is an explicit polynomial identity
\[
 N_G(X)=U(X,Y)G(X,Y)+V(X,Y)F(X,Y),
 \tag{NM}\label{eq:norm-membership}
\]
where \(U=(b\operatorname{adj}(M_G))_0\) and
\(V=-(h\operatorname{adj}(M_G))_0\).  Thus no ideal-membership
oracle is needed.  If \(N_G=c\ne0\) is constant, division of
\eqref{eq:norm-membership} by \(c\) gives the claimed identity
for \(1\).

Otherwise divide \(N_G\) by its nonzero rational leading
coefficient to obtain a monic \(q(X)\) of degree \(d\), retaining
the correspondingly scaled witnesses in \eqref{eq:norm-membership}.
Use the displayed monomial rectangle as the generating family, with
\(1\) its first member.  Multiplication by \(X\) stays inside
the rectangle unless \(i=d-1\); in that case the monic relation
\(q(X)\), multiplied by \(Y^j\), gives the action coefficients
and their ideal witnesses directly.

Multiplication by \(Y\) stays inside the rectangle unless
\(j=m-1\).  In that case write
\[
 -X^i f_t(X)=s_t(X)q(X)+r_t(X),\qquad \deg r_t<d
 \quad(0\le t<m)
\]
by bounded monic division in \(X\).  The exact identity
\[
 YX^iY^{m-1}-\sum_{t<m}r_t(X)Y^t
       =X^iF+\sum_{t<m}s_t(X)q(X)Y^t
\]
provides this boundary column.  Substitute the already retained ideal
witness for \(q\) on the right.  All remaining terms are explicit
multiples of \(F,G\), and the left coefficient combination lies
inside the rectangle.  These are precisely the required unit and
action identities.

The division in \(Y\) has its input-degree bound, the determinant
and adjugate are finite permutation sums, and each subsequent monic
division in \(X\) has its input-degree bound.  The constructor
performs these fixed loops over \(md\) monomials and retains the
identities at each step.  Its proof uses coefficientwise finite
identities as bounded invariants.  It is therefore primitive recursive
uniformly in the supplied polynomial codes and has the indicated
iPRA correctness proof strategy.
\end{proof}

The rational-root scan is usually smaller than the multiplication
certificate produced by this constructor.  Both have value: the scan
gives a direct implementation, while the constructor lets the same
general proof-producing finite-algebra interface certify the result.

\subsection{Existing weighted norms become effective section lists}

The previously supplied norm theorems treat, among other cases,
\[
 F(X,Y)=Y^m-cX^d+\text{terms of smaller weight},
 \quad \mathrm{wt}(X)=m,\quad\mathrm{wt}(Y)=d,
\]
where \(m\ge2\), \(d\ge1\), \(\gcd(m,d)=1\) and \(c\ne0\).
For a nonzero reduced polynomial \(G\), with \(\deg_YG<m\), those
theorems give an explicit nonzero leading coefficient for \(N_G\).
The new enumeration theorem uses precisely this nonvanishing output;
it need not reprove the sharper norm degree or fiber multiplicity bound.

Suppose a proper rational linear subspace in a finite monomial
evaluation space is supplied by a finite matrix.  Bounded rational
elimination produces a nonzero annihilator \(\lambda\); the
corresponding polynomial \(G_\lambda\) is a nonzero linear
combination of distinct reduced monomials.  The weighted theorem
certifies \(N_{G_\lambda}\ne0\), so
Theorem~\ref{thm:norm-complete-intersection} lists every rational
curve point on that section.  If only one annihilator is used, the
list may include points in the containing hyperplane but outside
the original subspace.  A final exact subspace-membership test
removes them.  Hence the section list is exact, not merely a cover
by candidates.

As a small example, take
\[
 F=Y^2-X^3+X-1,\qquad G=X^2-X.
\]
Since \(G\) has no \(Y\)-dependence, \(M_G=G I_2\), whence
\(N_G=(X^2-X)^2\).  Its rational-root list is \((0,1)\).
At either first coordinate, \(F_a(Y)=Y^2-1\), with roots
\(-1,1\).  The last filter retains exactly
\[
 (0,-1),\ (0,1),\ (1,-1),\ (1,1).
\]
This is the sharp four-point example already used for counting in
the older framework; here the same finite section data give a
certified exhaustive search.  The argument establishes a new
connection between existing packet interfaces, not priority for
the underlying rational-root or elimination methods.

\subsection{Global covers, tuple-local covers, and bounded coefficients}

Let \(S(a,b)\) be any fixed decidable condition on rational points
of \(F=0\), such as integrality.  Suppose actual polynomials
\(G_1,\ldots,G_s\) are supplied, each with nonzero norm, together
with a proof of the single global assertion
\[
 \forall a,b\quad
 F(a,b)=0\ \wedge\ S(a,b)
 \ \Longrightarrow\ \bigvee_{j=1}^{s}G_j(a,b)=0.
 \tag{GC}\label{eq:norm-global-cover}
\]
The cover assertion can be a supplied object-theory derivation;
merely checking its truth on an initial search box is insufficient.

\begin{corollary}[A supplied global cover gives all points]
\label{cor:norm-global-list}
Under \eqref{eq:norm-global-cover}, concatenate the lists
\(\mathcal I(F,G_j)\), remove duplicate pairs and retain the
pairs satisfying \(S\).  The result is a complete finite list of
all rational points satisfying \(F=0\) and \(S\).
The maximum of the bounds \(H(F,G_j)\), or one for an empty
cover, is a height bound.
\end{corollary}

\begin{proof}
Every retained pair is a curve point satisfying \(S\).  Conversely,
a point satisfying these conditions belongs to one supplied section
by the finite disjunction in \eqref{eq:norm-global-cover}.  The
corresponding intersection theorem puts it in that section's list.
Concatenation and deduplication preserve membership.  This is a
finite intuitionistic disjunction elimination, not a selection of
an infinite family of exceptional subspaces.  If the cover is
empty, its global assertion already proves that no such point exists.
\end{proof}

The Siegel route in the repository instead asks that \emph{each
supplied finite tuple} admit some cover by at most \(M\) proper
subspaces.  Its sections may change with the tuple, and their
coefficients may depend on the tuple's coordinates.  Applying the
present theorem to a cover of an already supplied tuple locates
only points of that cover; it gives no guarantee that all further
points lie there.  Even a sharp bound on the number of points in
each section does not repair this quantifier change.  Thus the
new enumeration theorem resolves the localization of explicit
sections, but does not supply the missing global approximation
cover or prove general Siegel or Faltings.

A useful sufficient strengthening can nevertheless be stated
precisely.  Fix a list of \(r\) distinct reduced monomials and a
positive integer \(C\).  There are at most
\((2C+1)^r-1\) nonzero integer coefficient vectors
\(\lambda\) with \(|\lambda_i|\le C\).  For each vector, compute
the norm of its monomial combination and retain those with nonzero
norm.  If a proof establishes that every point in \(S\) lies on
one of these retained sections, the preceding corollary computes
all such points.  In the weighted coprime setting, every nonzero
combination is automatically retained.  Consequently a
\emph{uniform coefficient-height bound} for the exceptional
annihilators would suffice, even if no preferred global cover were
selected.  Tuple-dependent coefficient bounds do not suffice.
This identifies a concrete missing quantitative input rather than
an unspecified need for choice.

For Mordell--Weil, none of these conclusions constructs a rational
point on every relevant descent covering.  Such covering curves are
typically positive-dimensional, and a finite quotient or nonzero
intersection norm becomes available only after an additional
independent equation is supplied.  Solving its resulting finite
intersection does not prove that some intersection contains a
representative of every required quotient class.  The rational
representative-list obligation therefore remains intact.

\subsection{Lean and object-iPRA obligations}

The definitions needed for a complete formalization are finite:
signed rational codes with normalization; univariate and bivariate
coefficient arrays; monic \(Y\)-reduction; a determinant fold;
\(\operatorname{RootBound}\), \(\operatorname{RootList}\), and
\(\operatorname{IntersectionList}\); and exact Boolean root and
subspace checks.  The recommended lemma order is:
\begin{enumerate}
\item prove gcd/Bezout cancellation, the denominator-cleared
polynomial identity, and rational-root bounds;
\item prove the bounded root scan's soundness and completeness;
\item prove monic reduction under evaluation and the adjugate
identity, then the norm-vanishing lemma;
\item combine the two root scans into complete intersection
enumeration, with an explicit maximum over first-coordinate candidates;
\item combine finitely many intersection lists under a supplied
derivation of the global cover formula.
\end{enumerate}
At the embedded level, each scan's invariant states that its output
contains exactly the successful candidates examined so far, as a
bounded Boolean equality.  Coefficient clearing, reduction and
determinant identities are equalities of finite codes or bounded
coefficient arrays.  Euclidean descent uses fixed numerical fuel.
Free point codes remain parameters throughout these inductions.
Only after constructing the arithmetic derivations does one use
the repository's native equivalence.  An ambient Lean theorem
about these arrays, or a type-checked declaration of the target,
does not on its own construct the required \texttt{Prf PRAAx} term.
